% ═══════════════════════════════════════════════════════════════════════════
% QR-CODE: Minitest Energieumwandlung & Erhaltungssatz (Beamer-Projektion)
% Physik Klasse 10 MR | Stunde 04
% ═══════════════════════════════════════════════════════════════════════════

\documentclass[a4paper,11pt]{article}

\usepackage[utf8]{inputenc}
\usepackage[T1]{fontenc}
\usepackage{helvet}
\renewcommand{\familydefault}{\sfdefault}

\usepackage[margin=20mm]{geometry}
\usepackage{tikz}
\usepackage{xcolor}
\usepackage{qrcode}
\usepackage{hyperref}

% ═══════════════════════════════════════════════════════════════════════════
% FARBEN
% ═══════════════════════════════════════════════════════════════════════════

\definecolor{physik}{HTML}{2c5aa0}
\colorlet{fachfarbe}{physik}

% ═══════════════════════════════════════════════════════════════════════════
% KONFIGURATION
% ═══════════════════════════════════════════════════════════════════════════

\newcommand{\mttitel}{Energieumwandlung \& Erhaltungssatz}
\newcommand{\mturl}{https://devbox42.github.io/sciencesim/physik/kl10-MR/04-energieumwandlung/MT-04-energieumwandlung.html}

% ═══════════════════════════════════════════════════════════════════════════
% DOKUMENT
% ═══════════════════════════════════════════════════════════════════════════

\pagestyle{empty}

\begin{document}

\begin{center}

% Titel
\vspace*{10mm}
{\fontsize{28}{34}\selectfont\bfseries\color{fachfarbe}Minitest: \mttitel}

\vspace{15mm}

% QR-Code mit Rahmen
\fcolorbox{fachfarbe}{white}{\qrcode[height=100mm]{\mturl}}

\vspace{12mm}

% URL - vollständig anklickbar
{\small\color{gray}\url{\mturl}}

\vspace{8mm}

% Hinweis
{\Large\color{gray}Scanne den QR-Code mit deinem Gerät}

\vspace{10mm}

% Info-Box
\fcolorbox{fachfarbe}{fachfarbe!10}{%
    \parbox{0.7\textwidth}{%
        \centering
        \vspace{3mm}
        {\large\bfseries 23 BE | 20 Minuten}\\[3mm]
        Wähle zuerst dein Niveau (MR/GY).\\[1mm]
        Nutze deine Platzkarte für die Abgabe.
        \vspace{3mm}
    }%
}

\end{center}

\end{document}
